\paragraph{边界--循环秩机制与曲率可和口径下的进一步强化}
本段在既有定义与定理之上，给出一组可直接审计的强化命题。
核心依赖链为：规范体积密度与界面面积的互补不等式、边界方程的循环秩计数、prime-register 外置预算下界，以及跨分辨率曲率概率误差条。

\leanverified{paper\_killo\_visible\_entropy\_density\_area\_control}
\begin{theorem}[可见熵密度受面积密度控制]\label{thm:killo-visible-entropy-density-area-control}
设
\[
f_n:\Omega_n=\{0,1\}^n\twoheadrightarrow X_n
\]
为满射粗粒化，\(U_n\sim\mathrm{Unif}(\Omega_n)\)，\(Y_n:=f_n(U_n)\)。
记
\[
a(f_n):=\frac{A(f_n)}{n2^{n-1}},
\qquad
g(f_n):=2^{-n}\log|\mathcal G(f_n)|.
\]
若 \(|X_n|/2^n\to 0\)（\(n\to\infty\)），则
\[
\frac{H(Y_n)}{n}\ \le\ (\log 2)\,a(f_n)+o(1).
\]
\end{theorem}
\begin{proof}
令 \(d_x:=|f_n^{-1}(x)|\)、\(p_x:=d_x/2^n\)（\(x\in X_n\)），则 \(H(Y_n)=H(p)\)。
由定理 \ref{thm:hypercube-gauge-volume-shannon-stirling-gap}，
\[
g(f_n)=n\log 2-1-H(Y_n)+\Xi_n,
\]
其中
\[
\Xi_n:=\frac1{2^{n+1}}\sum_{x\in X_n}\log(2\pi d_x)+R(f_n),\qquad
0\le R(f_n)\le \frac{|X_n|}{12\cdot 2^n}.
\]
又 \(d_x\le 2^n\)，故
\[
0\le \Xi_n
\le
\frac{|X_n|}{2^n}\!\left(\frac{n\log 2+\log(2\pi)}{2}+\frac1{12}\right)
=o(n),
\]
从而
\[
\frac{g(f_n)}{n}=\log 2-\frac{H(Y_n)}{n}+o(1).
\]
另一方面由定理 \ref{thm:spg-gauge-volume-area-complementarity}，
\[
\frac{g(f_n)}{n}+(\log 2)\,a(f_n)\ge \log 2-\frac1n.
\]
代入并整理即得
\[
\frac{H(Y_n)}{n}\le (\log 2)\,a(f_n)+o(1).
\]
证毕。
\end{proof}

\begin{corollary}[正信息率排斥常数外置记忆]\label{cor:killo-positive-entropy-rate-excludes-constant-external-memory}
在定理 \ref{thm:killo-visible-entropy-density-area-control} 的设定下，若存在常数 \(h>0\) 使
\[
\liminf_{n\to\infty}\frac{H(Y_n)}{n}\ge h,
\]
则
\[
\liminf_{n\to\infty}a(f_n)\ge \frac{h}{\log 2}>0.
\]
进一步，固定素数 \(p\) 与轴容量参数 \(E\ge 0\)。
若每个 \(n\) 上都要求存在注入
\[
\iota_n:\FF_p(\delta_n)\hookrightarrow \mathcal P_{k_n,E},
\qquad
\delta_n\in\mathrm{im}(\partial_p),
\]
则必有
\[
k_n\ \ge\ \frac{\log p}{\log(E+1)}
\Bigl(\frac{h}{\log 2}+o(1)\Bigr)\,n2^{n-1}.
\]
因而不可能存在 \(k_n=O(1)\) 的常数轴数外置单射化方案。
\end{corollary}
\begin{proof}
第一式由定理 \ref{thm:killo-visible-entropy-density-area-control} 直接给出。

令
\[
r(f_n):=A(f_n)-|X_n|+c(B_{f_n})
\]
（定理 \ref{thm:spg-boundary-affine-fiber-cycle-rank}）。
由 \(c(B_{f_n})\le |X_n|\) 与 \(|X_n|/2^n\to 0\) 得
\[
r(f_n)=A(f_n)+o(n2^{n-1})=(a(f_n)+o(1))\,n2^{n-1}.
\]
再由定理 \ref{thm:spg-prime-register-budget-lower-bound}，
\[
k_n\log(E+1)\ge r(f_n)\log p.
\]
结合 \(a(f_n)\ge h/\log 2+o(1)\) 得
\[
k_n\ge
\frac{\log p}{\log(E+1)}
\Bigl(\frac{h}{\log 2}+o(1)\Bigr)\,n2^{n-1}.
\]
因此 \(k_n\) 至少线性于 \(n2^{n-1}\)，不可能为 \(O(1)\)。
证毕。
\end{proof}

\leanverified{paper\_killo\_positive\_entropy\_rate\_excludes\_constant\_external\_memory}

\leanverified{paper\_killo\_fold\_approx\_address\_fano\_lower\_bound}
\begin{theorem}[Fold 近似地址化的 Fano 下界]\label{thm:killo-fold-approx-address-fano-lower-bound}
固定分辨率 \(m\)。
令 \(\mathbf A\sim\mathrm{Unif}(\Omega_m)\)，\(X:=\Fold_m(\mathbf A)\in X_m\)，并记
\[
d_m(x):=|\Fold_m^{-1}(x)|.
\]
取有限记录轴 \(r:\Omega_m\to R\) 与解码器 \(D:X_m\times R\to\Omega_m\)。
若
\[
\PP\!\bigl[D(X,r(\mathbf A))=\mathbf A\bigr]\ge 1-\delta,\qquad 0<\delta<1,
\]
则
\[
\log|R|\ \ge\ \EE\!\bigl[\log d_m(X)\bigr]-h(\delta)-\delta\,m\log 2,
\]
其中 \(h(\delta)=-\delta\log\delta-(1-\delta)\log(1-\delta)\)。
特别地，若 \(\EE[\log d_m(X)]\ge c\,m\)（某 \(c>0\)）且 \(|R|=O(1)\)，则 \(\delta\) 不能趋于 \(0\)。
\end{theorem}
\begin{proof}
令 \(\widehat{\mathbf A}:=D(X,r(\mathbf A))\)，则
\(
\PP[\widehat{\mathbf A}\neq \mathbf A]\le \delta
\)。
由 Fano 不等式，
\[
H(\mathbf A\mid X,r(\mathbf A))
\le
h(\delta)+\delta\log(|\Omega_m|-1)
\le
h(\delta)+\delta\,m\log 2.
\]
又
\[
H(\mathbf A\mid X)
=
H(\mathbf A\mid X,r(\mathbf A))+I(\mathbf A;r(\mathbf A)\mid X),
\]
且
\[
I(\mathbf A;r(\mathbf A)\mid X)\le H(r(\mathbf A)\mid X)\le \log|R|.
\]
故
\[
\log|R|
\ge
H(\mathbf A\mid X)-h(\delta)-\delta\,m\log 2.
\]
最后计算 \(H(\mathbf A\mid X)\)。
对任意 \(x\in X_m\)，在条件 \(X=x\) 下，\(\mathbf A\) 在纤维 \(\Fold_m^{-1}(x)\) 上均匀分布，故
\[
H(\mathbf A\mid X=x)=\log d_m(x).
\]
取期望得
\[
H(\mathbf A\mid X)=\EE[\log d_m(X)].
\]
代回即证。
\end{proof}

\begin{proposition}[曲率事件可和蕴含最终一致的逆极限提升]\label{prop:killo-curvature-summable-eventual-inverse-limit-lift}
\leanverified{paper\_killo\_curvature\_summable\_eventual\_inverse\_limit\_lift}
记 \(\omega^{(m)}:=\tau_{\infty\to m}(\omega_\infty)\)。
若
\[
\sum_{m\ge 1}\PP\!\left[\mathcal K_{m+1\to m}(\omega^{(m+1)})=1\right]<\infty,
\]
则几乎处处存在 \(m_0(\omega_\infty)\) 使得对所有 \(m\ge m_0\) 成立
\[
\pi_{m+1\to m}\!\bigl(\Fold_{m+1}(\omega^{(m+1)})\bigr)
=
\Fold_m(\omega^{(m)}).
\]
从而尾序列 \(\bigl(\Fold_m(\omega^{(m)})\bigr)_{m\ge m_0}\) 定义了唯一的逆极限元素
\[
x_\infty(\omega_\infty)\in X_\infty\simeq \varprojlim X_m.
\]
\end{proposition}
\begin{proof}
设事件
\(
A_m:=\{\mathcal K_{m+1\to m}(\omega^{(m+1)})=1\}
\)。
由第一 Borel--Cantelli 引理，\(\sum_m\PP(A_m)<\infty\Rightarrow \PP(A_m\ \mathrm{i.o.})=0\)。
故几乎处处仅有限多个 \(m\) 使 \(A_m\) 发生，即存在 \(m_0\) 使 \(m\ge m_0\) 时 \(\mathcal K_{m+1\to m}=0\)。
按定义 \(\mathcal K_{m+1\to m}=0\) 等价于
\[
\pi_{m+1\to m}(\Fold_{m+1}(\omega^{(m+1)}))=\Fold_m(\omega^{(m)}).
\]
这给出逆系统相容尾部，因此确定唯一逆极限点。
证毕。
\end{proof}

\begin{theorem}[曲率概率控制跨分辨率观测极限一致性]\label{thm:killo-curvature-probability-observable-limit-consistency}
\leanverified{paper\_killo\_curvature\_probability\_observable\_limit\_consistency}
在命题 \ref{prop:killo-curvature-summable-eventual-inverse-limit-lift} 的设定下，记
\[
\varepsilon_m:=\PP\!\left[\mathcal K_{m+1\to m}(\omega^{(m+1)})=1\right],\qquad
\sum_{m\ge 1}\varepsilon_m<\infty.
\]
取任意固定分辨率 \(m_\ast\) 与任意有界可观测量 \(f:X_{m_\ast}\to\RR\)。
定义
\[
Z_m:=\EE\!\left[f\!\left(\pi_{m\to m_\ast}(\Fold_m(\omega^{(m)}))\right)\right],\qquad m\ge m_\ast.
\]
则 \((Z_m)_{m\ge m_\ast}\) 收敛。
并且“先沿 \(\tau\) 再 Fold”与“先 Fold 再沿 \(\pi\)”两条路径在 \(m\to\infty\) 下给出同一极限。
\end{theorem}
\begin{proof}
令 \(f_m:=f\circ \pi_{m\to m_\ast}:X_m\to\RR\)。
由命题 \ref{prop:pom-gauge-curvature-shadow-bound}（取 \(\mu=\mathcal L(\omega^{(m+1)})\) 与可观测量 \(f_m\)）得
\[
\begin{aligned}
|Z_{m+1}-Z_m|
&=
\left|
\EE f_m\!\left(\pi_{m+1\to m}(\Fold_{m+1}(\omega^{(m+1)}))\right)
-\EE f_m\!\left(\Fold_m(\tau_{m+1\to m}(\omega^{(m+1)}))\right)
\right|\\
&\le 2\|f_m\|_\infty\,\varepsilon_m
\le 2\|f\|_\infty\,\varepsilon_m.
\end{aligned}
\]
于是对 \(N>M\ge m_\ast\)，
\[
|Z_N-Z_M|
\le
\sum_{m=M}^{N-1}|Z_{m+1}-Z_m|
\le
2\|f\|_\infty\sum_{m=M}^{N-1}\varepsilon_m.
\]
右端由可和性趋于 \(0\)，故 \((Z_m)\) 为 Cauchy 列并收敛。

再记
\[
Y_m:=\EE\!\left[f\!\left(\Fold_{m_\ast}(\tau_{m\to m_\ast}(\omega^{(m)}))\right)\right].
\]
对固定 \(m\ge m_\ast\)，定义中间量
\[
W^{(m)}_j:=
\EE\!\left[
f\!\left(\pi_{j\to m_\ast}\!\bigl(\Fold_j(\tau_{m\to j}(\omega^{(m)}))\bigr)\right)
\right],
\qquad
j=m_\ast,\dots,m.
\]
则 \(W^{(m)}_{m_\ast}=Y_m\)、\(W^{(m)}_m=Z_m\)。
对每个 \(j\in\{m_\ast,\dots,m-1\}\)，把命题 \ref{prop:pom-gauge-curvature-shadow-bound} 应用于可观测量 \(f_j=f\circ\pi_{j\to m_\ast}\) 与随机变量 \(\tau_{m\to j+1}(\omega^{(m)})\sim\mu_{j+1}\)，得
\[
\left|W^{(m)}_{j+1}-W^{(m)}_j\right|
\le 2\|f\|_\infty\,\varepsilon_j.
\]
沿 \(j\) 求和即得
\[
|Y_m-Z_m|
\le
2\|f\|_\infty\sum_{j=m_\ast}^{m-1}\varepsilon_j
\xrightarrow[m\to\infty]{}0.
\]
因此两条路径极限相同。
证毕。
\end{proof}

\begin{theorem}[单一水库纤维内的线性 VC 维下界]\label{thm:killo-single-reservoir-fiber-linear-vc-lower-bound}
\leanverified{paper\_killo\_single\_reservoir\_fiber\_linear\_vc\_lower\_bound}
取 \(m=4k\)，并沿用定义 \ref{def:fold-block-reservoir-code} 的
\[
\rho_k=(1000)^k,\qquad
\mathrm{Enc}_k:\{0,1\}^k\to \Omega_m.
\]
令纤维
\[
F:=\Fold_m^{-1}(\rho_k).
\]
定义概念类
\[
\mathcal C:=\{C_T:\ T\subseteq [k]\},
\]
其中
\[
C_T:=\left\{\omega\in F:\ \forall j\in T,\ \omega|_{\text{第 }j\text{ 块}}=0110\right\}.
\]
则
\[
\mathrm{VCdim}(\mathcal C|_F)\ \ge\ k=\frac{m}{4}.
\]
\end{theorem}
\begin{proof}
由定理 \ref{thm:fold-block-reservoir-encoding-constant-radius}，\(\mathrm{Enc}_k\) 为单射且
\[
\Fold_m(\mathrm{Enc}_k(w))=\rho_k\qquad(\forall w\in\{0,1\}^k),
\]
故 \(\mathrm{Enc}_k(\{0,1\}^k)\subseteq F\)。

对 \(i\in[k]\)，令 \(p_i\in\{0,1\}^k\) 满足 \(p_i(i)=0\) 且 \(p_i(j)=1\)（\(j\neq i\)），并置
\[
P:=\{\mathrm{Enc}_k(p_i):\ i=1,\dots,k\}\subseteq F.
\]
任取 \(S\subseteq[k]\)，令 \(T:=[k]\setminus S\)。
若 \(i\in S\)，则对每个 \(j\in T\) 有 \(p_i(j)=1\)，故 \(\mathrm{Enc}_k(p_i)\) 的第 \(j\) 块为 \(0110\)，从而 \(\mathrm{Enc}_k(p_i)\in C_T\)。
若 \(i\notin S\)，则 \(i\in T\)，而 \(p_i(i)=0\)，故第 \(i\) 块为 \(1000\neq 0110\)，从而 \(\mathrm{Enc}_k(p_i)\notin C_T\)。
于是
\[
C_T\cap P=\{\mathrm{Enc}_k(p_i):\ i\in S\}.
\]
即 \(P\) 被 \(\mathcal C\) 打碎。
按定义 \ref{def:conclusion-vc-dim}，得到
\[
\mathrm{VCdim}(\mathcal C|_F)\ge |P|=k.
\]
证毕。
\end{proof}

\endinput
